\chapter[Conclusion]{Conclusion}
\label{Chap:Conclusion}

This thesis examined how kinematically narrow and broad \HI\ are distributed and associated with star formation in 16 nearby galaxies. We used ROHSA to decompose the \HI\ profiles and compared the separated maps with CO-based H$_2$ and SFR maps. We use the kinematically cold and warm \HI\ components as proxies for the CNM and WNM.

Both atomic components extend across much of the mapped \HI\ discs. Across the sample, the broad component contributes a median of about 55\% of the total \HI\ mass, while the kinematically cold component accounts for the remaining 45\%. The kinematically cold component has a more filamentary appearance, but its fractional contribution is not consistently higher in stellar spiral arms. We also find no clear evidence that narrow \HI\ is generally more concentrated near CO emission than broad \HI. This may be explained by how atomic gas behaves during the atomic-to-molecular transition where atomic gas remains in shielding layers around molecular interiors while additional gas becomes molecular. Hydrogen also no longer contributes to the \HI\ emission once it is converted into H$_2$. Additionally, H$_2$ can survive where CO is dissociated, so our CO-faint regions are not necessarily free of H$_2$.

We examined the Kennicutt--Schmidt relations for the different gas components. Their behaviour depends on the galaxy population. In the combined dwarf sample, narrow \HI\ correlates more strongly with the SFR than total or broad \HI. In spirals, broad \HI\ follows the SFR more closely than narrow \HI. One possible explanation for the stronger broad-component relation in spirals is that turbulence broadens the emission of cold gas near star-forming regions, allowing some of it to be assigned to the broad component.

As already established in the literature, the molecular component remains more closely associated with the SFR in most spirals and follows an approximately linear Kennicutt--Schmidt relation \citep{Fbigiel2008}.

The SFE of both the narrow and broad \HI\ components declines towards the outer disc. In dwarfs, the narrow \HI\ component has a higher SFE than the broad component across the disc. In spirals, narrow \HI\ has the higher SFE in the inner disc, while the efficiencies of the narrow and broad components become more similar farther out. One possible explanation for this convergence is that turbulent broadening becomes weaker towards the outer disc, consistent with the radial decline in \HI\ line widths reported by \citet{Tamburro2009}. A larger fraction of the cold gas may then contribute to the narrow component rather than the broad component, bringing their relative contributions and measured SFEs closer together.

Our data do not allow us to locate where the \HI-to-H$_2$ transition occurs, since molecular hydrogen can be present without detectable CO emission. Extending the analysis to a larger sample and incorporating data from more recent surveys with improved spatial and spectral resolution, independent constraints on the thermal phases and molecular-gas measurements, accounting for metallicity or other parameters would help clarify the role of cold atomic gas. A larger dwarf sample would also help establish whether the stronger combined narrow-\HI\ relation is a general property of these galaxies.